\chapter{Заключение}

В крайния си вид проектът представлява работещо .NET приложение, което покрива съществените изисквания на практическото задание и в същото време остава смислено от инженерна гледна точка. Системата използва MVVM, поддържа две различни View технологии, използва споделени ViewModel-и, работи с SQLite и PostgreSQL чрез общ модел и общ бизнес слой и включва преизползваеми функционалности за филтриране и визуализиране на списъци от обекти.

Най-важният архитектурен извод от разработката е, че правилното разделение по слоеве дава по-добър резултат от механично разпределение на файловете по шаблонни папки. Чрез такова разделение проектът постигна споделена приложна логика, технологично независими ViewModel-и, конфигурационно превключване между две СУБД и преизползваеми абстракции за потребителски интерфейс, без да жертва разбираемостта на реализацията.

От практическа гледна точка разработката показа и още нещо важно. За добър курсов проект не е достатъчно кодът да бъде архитектурно коректен. Необходимо е системата да бъде стабилна в реална среда, демонстрационно ясна и функционално завършена. Именно поради това в проекта беше обърнато внимание не само на модела и слоя за съхранение, но и на фоновата обработка, обработката на видео, синхронизацията между двата клиента и подобренията в Avalonia интерфейса.

Макар текущият вариант вече да е достатъчно завършен за целите на заданието, проектът има потенциал и за бъдещо развитие. Възможни посоки са по-богато визуализиране на резултатите, въвеждане на миграции вместо \texttt{EnsureCreated}, по-пълно изравняване на функционалността между двата клиента и допълнителни аналитични функции върху резултатите от детекцията.
